\subsection{Two ways to turn a quantum Hamiltonian into classical
input}\label{quant:sec-1}

Let \(\mathcal H\) be an \(n\)-qubit Hamiltonian with a fixed
computational basis. Two possible classical input measures are

\[
\begin{aligned}
\mu_{\beta,\mathcal H}^{\rm meas}(x)
&=\frac{\langle x|e^{-\beta\mathcal H}|x\rangle}{\operatorname{Tr}(e^{-\beta\mathcal H})},\\
\mu_{\beta,\mathcal H}^{\rm diag}(x)
&=\frac{e^{-\beta\langle x|\mathcal H|x\rangle}}
{\sum_y e^{-\beta\langle y|\mathcal H|y\rangle}}.
\end{aligned}
\]

The first measures a quantum thermal state in the computational basis.
The second discards off-diagonal entries before taking a classical Gibbs
law. They agree for diagonal Hamiltonians but generally differ:
dephasing and exponentiating do not commute.

For the EPR interaction \(K=-|\psi^+\rangle\langle\psi^+|\), the
unnormalized measured thermal weights on \(00,01,10,11\) are

\[
1,\quad(1+e^\beta)/2,\quad(1+e^\beta)/2,\quad1.
\]

Exponentiating diagonal energies instead gives

\[
1,\quad e^{\beta/2},\quad e^{\beta/2},\quad1.
\]

These differ for every \(\beta>0\). An AA problem with quantum
interactions must choose a model. The counterexample below works for
both.
